\par\noindent\textbf{(1)} 取球内像电荷q'=−qa/d放于z'=a²/d。球面r=a处两源距离满足|r−dẑ|=(d/a)|r−z'ẑ|，故两势相消：Φ=[q/|r−dẑ|+q'/|r−z'ẑ|]/(4πε0)=0。外部Poisson方程、边界及无穷远衰减唯一，因而外域电势由两源严格表示。
\par\noindent\textbf{(2)} 真实电荷受像电荷力F\_z=(1/(4πε0))qq'/(d−z')²=−q²ad/[4πε0(d²−a²)²]，沿−z方向指向导体球。像电荷不是真正的内部分布电荷，力只用于外部场效应，不能把像源位置当成实际载流粒子。
\par\noindent\textbf{(3)} 包围球的高斯面但不包围外部点电荷，总感应电荷Q\_ind=q'=−qa/d；由外域镜像场远区看，点电荷与感应球总电荷为q−qa/d。接地球可与地线交换电荷，不能用电荷守恒强制Q\_ind=0。
\par\noindent\textbf{(4)} d≫a时Q\_ind≈−qa/d、|F|≈q²a/(4πε0 d³)，符合诱导极化弱响应；d→a+时|F|∼q²/[16πε0(d−a)²]，近平面像力。孤立中性球还须加入球心像电荷+qa/d以恢复总感应零，力不再与接地情况相同。
